\section[Realization in the Pre-Socratics]{Realization in the Pre-Socratics: Rene Guenon and Peter Kingsley, Some Observations}

In the final chapter of Rene Guenon's \textit{Metaphysical Principles of Calculus} he concludes his exposition dedicating the chapter to the paradoxes of the Eleatic, Zeno, disciple of Parmenides. Guenon's position is that Zeno's paradoxes are not examples of an emerging scientific ``rationalism" as the academic position mostly holds, but a metaphysical and esoteric teaching.

Writing with the credentials of an academic, but voice and authority of the metaphysician, Peter Kingsley throughout his works such as \textit{Reality}, \textit{In the Dark Places of Wisdom}, and \textit{Ancient Philosophy, Mystery and Magic}, takes a similar approach, insisting that not only have the academics gotten Parmenides, Empodecles, Zeno, and related pre-Socratics wrong from ``back to front", not just that what they were not were dry proto-scientists—but as links in the transmission of an esotericism, an esotericism moreover that lays at the very roots of Western civilization, and that their doctrines can never really be subject to any scholarly ``debate" —rather, the only way to understand them, is that they ``must be lived".

If the goal that Guenon's works point toward might need summarizing, this perhaps could be done with just the word ``Realization", or perhaps attainment of the ``Supreme Identity" —if we elaborate, as a starting point maybe also taking ``man as envisaged in his totality"—in other words, quite simply, knowledge of the Real. Everything else in between—be it castes and civilizations, religions, initiations, social organization, the arts, crafts, and sciences, or teachings about ``cyclicity" (microcosmic or macrocosmic), in the traditional worldview, are all but images of the Real, and ultimately just the sign posts pointing back to it. And, ``outside" of the traditional world, where these things are unknown, there is not so much of an affirmative opposition to such things, as a radical absence, as Frithjof Schuon remarks,

\begin{quotex}
The spirits of evil, the demons, are inverted shadows—inverted in the direction of nothingness, which itself is in the nonexistent—of the Names of God; the ``fall of the angels" indicates the cosmic manifestation of the principles of remotion, inversion, privation, and negation, as well as compression and volatilization. Does this mean that the ideas of ugliness or vice are themselves a participation in ugliness or vice, as is the case with positive ideas? Clearly not, for the definition of vice is derived from virtue; goodness is the measure of evil; it is not because of stupidity that we know stupidity but because of intelligence, which makes recognition of this privation possible. \flright{\textit{Logic and Transcendence}, 19}

\end{quotex}
Quite akin to Guenon, Kingsley and his understanding of Parmenides and Empodecles concern much the same—we are asleep and unconscious, groping through a dream and illusion, a grand cosmic deception, and only by awaking and acquiring consciousness (rather realizing what we already possess) can we experience ``Reality".  The Real must be experienced, no other way is open —the ``Supreme Identity" must be ``realized", or it is but a virtuality. Since there is no way to really know, or cut through the ``deception" without ``realization", to ``describe" Reality, is really the same thing as transmitting the ``Initiatic Secret"—this is the Ineffable Word.

Still, Empodecles offers some framework or reference points: the nature of the four elements (or ``roots"), their joining and separation—and the point that their joining and separation weave the fabric of Reality —perpetually… cosmic cycle, after cosmic cycle, after cosmic cycle— which has some pretty important implications concerning our liberation, ``final deliverance"…or knowledge of Reality.

As is well known, Rene Guenon also placed some emphasis on cosmic ``cyclic" doctrines and symbolism, particularly the Indian, although often exploring relationships with the Persian, Chaldean, Islamic, and others. For Guenon, these cyclic doctrines were valuable in helping to assess the outward conditions of the world where it currently ``is", the impact this has on man and his civilizations, with the ever developing ``solidification" of the world being explained as a sort of growing ``distance" between the world and the principle of its origin —a ``distance" which like the ``fall" explained by Schuon earlier, has no ``positive existence", but is privation. In the macrocosmic arena, what is cyclic ``devolution" resulting in ``solidification", is in the individual being the degree of their unconsciousness —their inability to know Reality.

Yet, as Guenon and others have said, the ``fall" is never total (for again, a total fall would be another negation and nothingness). For the individual, the fall stops, and Reality becomes known, when plurality is resumed within unity, its self-sufficient principle, as Kingley suggests:

\begin{quotex}
When we are absorbed by everything happening all around us and excited with so many changes taking place… there is always movement. But when we start to glimpse the endless repetition behind our restless thirsting and questing for change, this is to catch sight of the motionlessness at the heart of all motion. For what can be so hard to appreciate is that motionlessness is not some different reality, some other level of existence, separate from the world of motion. On the contrary: the more we try to discover stillness by leaving movement behind, the more movement we end up creating through the very act of running away from motion. And just as Parmenides knew that the easiest way to experience the stillness seemingly denied by this world of the senses is by using our senses to the full, so Empodecles knew how to look for it in the place we would least expect to find it. To observe life with complete awareness is to see that, once movement has repeated itself long enough, the motion turns to motionlessness. \flright{\textsc{Kingsley}, \textit{Reality}, 472}

\end{quotex}
What in the individual constitutes the expansion of ``total awareness", via a ``stillness" bringing about and the ``end of motion", is in the macrocosm analogous to the end of a cycle, Guenon explains:

\begin{quotex}
If carried to its extreme limit the contraction of time would in the end reduce it to a single instant, and then duration would really have ceased to exist for it is evident that there can no longer be any succession within the instant….As soon as succession has come to an end, or in symbolic terms `the wheel has ceased to turn', all that exists can not but be in perfect simultaneity (``Reign of Quantity and Sign of the Times", 160).

\end{quotex}
However, as Empodecles, and related cyclic doctrines inform, as even observation of natural cycles such as the daily, lunar, and annual display, one terminal point in a certain cycle is precisely the initiating point in one to follow, Guenon continues:

\begin{quotex}
The end of a cycle….is really only the end of the corporeal world itself in quite a relative sense….in reality the corporeal world is not annihilated but `transmuted' and immediately receives a new existence, because, beyond the `stopping-point' corresponding to the unique instant at which time is no more, `the wheel begins to turn again for the accomplishment of another cycle'. \flright{\textit{Reign of Quantity,} 162}

\end{quotex}
Far reaching is the scope of these workings of reality, Kingsley states, included within and extending to all things, individual, formal, and informal, alike —or, as Guenon might say, in every state of Being (or in the geometrical symbolism of the cross, in every degree of individual manifestation, and every degree of Universal manifestation—that is, every ``world").

As Kingsley soon notes though, with everything therein considered, a ramification that becomes obvious is a certain lack of freedom, because ``every soul is dragged out of heaven once more at the beginning of a new cycle, as it most definitely will be, the illusion will start all over again (Reality, 469)". In fact, to the ire of modern man, as Kingsley explains this tradition —no free will exists at all.

For Guenon, freedom is but a ``possibility". As an ``absence of constraint", freedom can enter no limitation, and consequently can only proceed from non-duality—thus, freedom begins as a metaphysical possibility—for the individual, freedom is but relative, because ``the Unity of Being is the principle of freedom in particular beings as well as in Universal Being, a being will be free to the extent it participates in this unity (\textit{Multiple States of Being}, 92)". This metaphysical unity Guenon mentions individual beings participating in, might then be seen as akin to the ``stillness" of Parmenides and Empodecles, Kingsley mentions, with his dismissal of free will in mind:

\begin{quotex}
There is a cosmic law….which states that whenever you come to the point of total pointlessness some other point is lying behind it…But you also know that the return of the soul to freedom is just as transitory; just as illusory….But inwardly its all over. The drama has no substance. To find that stillness means neither being drawn into the illusions nor trying to escape them. It means that, even when trapped we are free. And even when apparently incarnated we know what it is, at any moment, not to be incarnated. For we are not only where we are. We are exactly where we will always be (473).

\end{quotex}
The cultivation of the ``stillness" producing this state involves, as seen earlier, achieving a ``total awareness" that ``shrinks" as it were all motion into motionlessness, or all quantity (plurality) into pure quality (unity) within the individual. Empodecles enlists a particular method:

\begin{quotex}
After the certainty that Love is good and Strife is bad comes the conviction that Love is bad, Strife is good. But after this certainty, too, has come and gone we are left with the growing sense that Empodecles is concerned with something far subtler than any fixed identification of either as simply good or bad. It all comes down to the need to discover, in each single instant, what happens to be good or bad for that particular moment. And only one factor can help us now. This is the razor sharp awareness, the totally focused sensitivity to the present moment, known to the Greeks as metis….metis is as silently meaningful to Parmenides as it is to him (Empodecles)….it lies so close to the core of both of their teachings. \flright{\textit{Reality}, 455}

\end{quotex}
This ``metis" then, which engenders ``kairos" —the ``moment", the ``eternal present"— the metaphysical freedom that is Guenon's ``Unity of Being", is as Kingsley translates Empodecle's poem, ``no mere mortal metis"; perhaps some would compare this with the contrast between dianoia and doxa verses theoria and noesis—but in any event, it is a divine quality, for:

\begin{quotex}
Love is divine. Strife is divine. Each one of the four elements or roots is divine. So whatever you see or hear is divine. And whatever you see or hear with is divine. This means that when you manage to separate your awareness into one pure element—not by thinking about it, simply by doing it—you are pure divinity….not to perceive, but to perceive…perceiving—to watch the perceptive process itself….Nothing is to be left out. Not the slightest preference is to be shown to one sense as opposed to any other. And this choiceless, all embracing awareness can only happen in one moment: right now (348, 511).

\end{quotex}
The motion made motioneless in the eternal now through divine perception, metis, perceiving Itself, shares another kinship in traditional symbolism–the ``wu wei" of the Taoist ``Transcendent Man" described by Guenon, one who has reached the goal, and attained ``metaphysical freedom" in actuality. Thus, perhaps aprapos, we conclude these very casual observations with Guenon's description of that state in his ``Great Triad", as translated by Mr. Kingsley:

\begin{quotex}
Heaven is as it were made manifest to them through him. Because his action, or rather his influence, is `central', it imitates—and from the point of view of the world of humanity it also `incarnates' the `Activity of Heaven'. This influence is `actionless', which means that it does not involve any external activity. The `One and only Man' exercises his role as `unmoved mover' from his position at the center. He controls everything without intervening in anything (127).

\end{quotex}


\flrightit{Posted on 2013-01-30 by Frater M }
